\begin{example}
\label{editorial-source-example-idempotent-matrices}
For another example, consider
$R = k[\{t_{ij}\}_{i, j = 1}^{n}]/\mathfrak a$, where $\mathfrak a$ is
the ideal generated by the entries of the product matrix $T^2-T$,
$T = (t_{ij})$. From linear algebra, we know that under the
$GL(n, k)$-action defined by $g, T \mapsto gTg^{-1}$, $T$ is
classified by its rank and each $T$ is conjugate to some
$\text{diag}(1, \ldots, 1, 0, \ldots, 0)$, which has $r$ 1's and $n-r$ 0's.
Every idempotent matrix with entries in $k$ belongs to exactly one
of these rank orbits. Indeed each vector has the decomposition
$v=Tv+(v-Tv)$ into $\operatorname{Im}(T)$ and $\ker(T)$;
their intersection is zero since $Tw=w$ on the image and $Tw=0$
on the kernel. Bases of these two subspaces give the stated
conjugacy, and conjugation preserves rank. To deduce the irreducible
components of $\operatorname{Spec}(R)$, including its nonclosed
points, we give a coordinate-ring proof below. It proves that the
rank loci are disjoint open and closed irreducible components over
every field $k$. To separate ranks in characteristic zero one may use
$f(t_{ij}) = trace(T) = \sum_{i = 1}^nt_{ii}$. In positive
characteristic cases, things are slightly more tricky since we
might have $trace(T) = 0$ even if $T \neq 0$. For instance, if $k$ has characteristic $3$,
$$
trace\left(
\begin{matrix}
1 & & \\
& 1 & \\
& & 1
\end{matrix}
\right) = 3 = 0
$$
For every $0\leq j\leq n$ put
$e_j(T)=\operatorname{tr}(\bigwedge^jT)$, with $e_0(T)=1$.
On a rank-$r$ idempotent this equals $\binom rj$ in $k$:
in a basis of the image followed by the kernel, the exterior map
has eigenvalue $1$ on the $j$-fold wedges of image basis vectors
and $0$ on all other basis wedges. The exterior basis change is
invertible, so it preserves the trace of the original exterior map.
Thus if $r<s$, $e_s$ is zero on rank $r$ and one on rank $s$.
The complete tuple $(e_1,\ldots,e_n)$ separates all ranks in every
characteristic. A single chosen entry need not do so: in
characteristic $3$, ranks $3$ and $4$ both give $e_3=1$, since
$\binom33=1$ and $\binom43=4=1$ in $k$.
The displayed function $e_3$ isolates rank $3$ when $n=3$;
it does not isolate that rank for arbitrary $n$.
\medskip\noindent
We now prove the component assertions in the original ring $R$,
without assuming that $k$ is algebraically closed or that $R$ is
reduced. Keep the original matrix $T$ and define its coefficients
$d_0,\ldots,d_n\in R$ by
$$
F(z)=\det(I_n-T+zT)=\sum_{r=0}^n d_rz^r.
$$
The identity $T^2=T$ gives, in matrices over $R[x,y]$,
$$
(I_n-T+xT)(I_n-T+yT)=I_n-T+xyT.
$$
Hence $F(x)F(y)=F(xy)$. Comparing the coefficient of each
$x^ry^s$ proves $d_r^2=d_r$ and $d_rd_s=0$ for $r\ne s$;
evaluation at $z=1$ gives $\sum_r d_r=1$.
Put $R_r=R/(1-d_r)$ and $X_r=\operatorname{Spec}(R_r)$.
Inverting the idempotent $d_r$ is the same as imposing $d_r=1$,
so $X_r=D(d_r)=V(1-d_r)$ in $\operatorname{Spec}(R)$.
In particular these loci are open and closed. The exact ring maps
$$
R\longrightarrow\prod_{r=0}^n R_r,
\qquad a\longmapsto(a\bmod(1-d_r))_r,
\qquad
(\overline{a_r})_r\longmapsto\sum_{r=0}^n d_ra_r
$$
are inverse: the second expression does not depend on lifts since
$d_r(1-d_r)=0$, and orthogonality and the sum identity verify both
composites on every element.

For any prime $\mathfrak p\subset R$, apply the image/kernel
decomposition above to the original matrix over $\kappa(\mathfrak p)$.
If its rank is $r$, its determinant polynomial is exactly $z^r$.
Thus $d_s$ has residue $1$ for $s=r$ and $0$ for $s\ne r$.
It follows that $X_r$ is precisely the residue-field rank-$r$
locus, including all nonclosed points. Each $X_r$ is nonempty:
evaluation at $P_r=\operatorname{diag}(1,\ldots,1,0,\ldots,0)$,
with exactly $r$ ones, is a $k$-point.

For irreducibility take
$A=k[g_{ij}:1\leq i,j\leq n][\det(G)^{-1}]$,
where $G=(g_{ij})$. This is an integral domain. The formula
$T\mapsto GP_rG^{-1}$ gives the coordinate-ring homomorphism
$$
R_r\longrightarrow A,\qquad
t_{ij}\longmapsto\sum_{a=1}^r g_{ia}(G^{-1})_{aj}.
$$
The inverse entries are in $A$ by the adjugate formula; the
substituted matrix is idempotent and has determinant polynomial
$z^r$, so the displayed homomorphism is well defined. Its spectrum
map $\operatorname{Spec}(A)\to X_r$ is surjective. In fact for
any $\mathfrak p\in X_r$, bases of the image and kernel over
$\kappa(\mathfrak p)$ provide an invertible matrix $G_{\mathfrak p}$
with $T_{\mathfrak p}=G_{\mathfrak p}P_rG_{\mathfrak p}^{-1}$.
Evaluation $A\to\kappa(\mathfrak p)$ at these entries has a
prime kernel whose contraction to $R_r$ is the given prime,
because its composite sends every original $t_{ij}$ to its
residue at $\mathfrak p$.
The continuous image of an irreducible space is irreducible:
a union of two relatively closed proper subsets would pull back
to two closed subsets covering the source, and surjectivity makes
both proper. Therefore each $X_r$ is irreducible. An irreducible
subset cannot meet two members of the disjoint open-and-closed
cover $(X_r)$, since that would separate it into two nonempty
relatively closed subsets. The same separation excludes a connected
subset from meeting two members. Each $X_r$ is itself irreducible
and connected. Hence the irreducible components and the connected
components of $\operatorname{Spec}(R)$ are exactly $X_0,\ldots,X_n$.

The same original coordinates give useful explicit charts.
For a subset $I\subseteq\{1,\ldots,n\}$ let $P_I$ be the diagonal
matrix with entry $1$ at indices in $I$ and $0$ elsewhere, and set
$$
Q_I=TP_I+(I_n-T)(I_n-P_I),\qquad q_I=\det Q_I.
$$
Multilinearity in the columns of $F(z)$, retaining their original
order, gives
$$
d_r=\sum_{|I|=r}q_I:
$$
the term indexed by $I$ chooses the column $Te_j$ for $j\in I$
and $(I_n-T)e_j$ for $j\notin I$.
Thus $X_r$ is covered by the opens $D(q_I)$ with $|I|=r$.
To check that each such open lies in $X_r$, compute
$TQ_I=Q_IP_I$. After inverting $q_I$ this gives
$T=Q_IP_IQ_I^{-1}$, so $F(z)=z^r$ there.

Write $K=Q_I^{-1}$ in this localization. From $KT=P_IK$,
the definition of $Q_I$ gives the exact identity
$$
I_n=KQ_I=P_IKP_I+(I_n-P_I)K(I_n-P_I).
$$
Consequently the two diagonal blocks of $K$, indexed by $I$ and
its complement without changing the original row and column labels,
are identity matrices. Introduce variables $u_{ab}$ for the ordered
pairs $(a,b)$ with exactly one of $a,b$ in $I$. Let $K_I(u)$ have
these off-diagonal entries and the two identity diagonal blocks.
Put $h_I(u)=\det K_I(u)$ and
$$
B_I=k[u_{ab}: \text{exactly one of }a,b\text{ belongs to }I]
[h_I(u)^{-1}].
$$
There are $2r(n-r)$ such variables and $h_I(0)=1$.
The matrix $T'=K_I(u)^{-1}P_IK_I(u)$ is idempotent, and
$$
\begin{aligned}
Q_I(T')
&=K_I(u)^{-1}\bigl(P_IK_I(u)P_I
 +(I_n-P_I)K_I(u)(I_n-P_I)\bigr)\\
&=K_I(u)^{-1}.
\end{aligned}
$$
Thus $q_I(T')=h_I(u)^{-1}$ is invertible. Sending each original
$t_{ab}$ to the corresponding entry of $T'$ defines
$R[q_I^{-1}]\to B_I$. Conversely, sending each $u_{ab}$ to the
corresponding entry of $Q_I^{-1}$ defines $B_I\to R[q_I^{-1}]$:
the diagonal blocks agree by the identity just proved, and
$h_I$ maps to $q_I^{-1}$. These homomorphisms are inverse.
On the original matrix the composite is $Q_IP_IQ_I^{-1}=T$;
on the matrix of chart variables it is $Q_I(T')^{-1}=K_I(u)$.
The inverted determinants have the corresponding inverse images.
We have therefore proved the exact chart isomorphism
$$
R[q_I^{-1}]\cong B_I.
$$

In particular every chart ring is an integral domain and is reduced.
These charts prove that $R$ itself is reduced. If $a\in R$ is
nilpotent, its image is zero in each chart, so for each $I$ some
positive power $q_I^{N_I}$ annihilates $a$. The finitely many
$q_I$ cover the spectrum, since their sum over all $I$ is
$\sum_r d_r=1$. The powers $q_I^{N_I}$ also generate the unit
ideal: a maximal ideal containing them would contain every $q_I$,
contrary to that sum identity. Taking a linear combination equal
to $1$ proves $a=0$.
Each factor $R_r$ is reduced and has irreducible spectrum, hence is
a domain: if $ab=0$ and neither factor were zero, reducedness would
make both $D(a)$ and $D(b)$ nonempty, whereas their intersection is
$D(ab)=\varnothing$, contradicting irreducibility.
All constructions commute with every field extension $L/k$:
$R\otimes_k L=L[t_{ij}]/((T^2-T)_{ij})$, with the same coefficient
polynomials $d_r$ and chart maps. Repeating the argument over $L$
shows that each $R_r\otimes_k L$ is a domain. Thus each component
is geometrically integral, in the sense of being integral after
every extension of the original base field.

We can also identify its defining ideal exactly. Let $J_\ell(T)$
be the ideal generated by the $\ell\times\ell$ minors, with
$J_0(T)=R$ and $J_{n+1}(T)=0$. The diagonal entries of
$\bigwedge^jT$ in its indexed exterior basis are the principal
$j\times j$ minors. Expanding the determinant therefore gives
$$
\det(I_n+uT)=\sum_{j=0}^n e_j(T)u^j,
\qquad
d_s=\sum_{j=s}^n(-1)^{j-s}\binom js e_j(T).
$$
For $s\geq\ell$, every minor appearing in the latter sum belongs
to $J_\ell(T)$, by Laplace expansion along any fixed $\ell$ of
its rows. Hence $(d_\ell,\ldots,d_n)\subseteq J_\ell(T)$.
Conversely, on a chart of $X_s$ with $s<\ell$, the identity
$T=Q_IP_IQ_I^{-1}$ gives
$$
\bigwedge^\ell T
=(\bigwedge^\ell Q_I)(\bigwedge^\ell P_I)
 (\bigwedge^\ell Q_I^{-1})=0,
$$
since $P_I$ has only $s$ nonzero diagonal entries.
All $\ell\times\ell$ minors therefore vanish on each chart of
$X_s$, hence in $R_s$ by the same finite-cover annihilator argument.
In the proved product decomposition of $R$, the ideal of elements
whose coordinates for $s<\ell$ are zero is exactly
$(d_\ell,\ldots,d_n)$. This proves
$$
J_\ell(T)=(d_\ell,\ldots,d_n)
\quad(0\leq\ell\leq n+1).
$$
For the complementary idempotent matrix the Laurent-polynomial
identity
$$
\det(T+z(I_n-T))
=z^n\det(I_n-T+z^{-1}T)
$$
shows that $d_s(I_n-T)=d_{n-s}(T)$. Applying the minor formula
to both matrices gives the exact component ideal
$$
(1-d_r)=J_{r+1}(T)+J_{n-r+1}(I_n-T).
$$
Indeed the right side is generated by all $d_s$ with $s\ne r$;
their sum is $1-d_r$ and each satisfies $d_s(1-d_r)=d_s$.
These are also precisely the minimal prime ideals of $R$, since
each quotient $R_r$ is a domain and the $X_r$ are all components.
Finally $F(1+u)=\det(I_n+uT)$ gives
$$
e_j(T)=\sum_{r=0}^n\binom rj d_r,
$$
an equality in the original ring, strengthening the residue-field
trace calculation. Every integer coefficient is mapped to $k$;
no characteristic restriction has been introduced. If $n=0$,
empty determinants are $1$, $R=k$, $d_0=1$, and the same statements
hold with their empty matrices and lists.
\end{example}